% !TeX root = ../MADRE-AgenticSystem.tex

\madreChapter{Safety And Evolution Policy}

\section{Runtime Safety}

\Madre is intended to run useful delayed work while the user is not constantly present. That requirement makes safety an architectural property, not a documentation afterthought. The system must be able to explain what it did, why it did it, what data it used, what boundary allowed it, and how the result can be corrected or rolled back when the action is reversible.

\subsection{Policy Before Action}

The kernel must not expose executable behavior merely because a model can request it. Internal actions, external calls, knowledge-boundary changes, workflow execution, and learning promotion need an explicit \texttt{PolicyDecision} before execution. The safe path should be the normal path, and safe autonomous improvement should not require constant user interruption. Unsafe paths should be technically isolated, traceable, recoverable where possible, and visible to the responsible developer or operator through \texttt{RuntimeJournal}.

Typed runtime-defined internal actions represent the canonical bounded action form. They may be allowed only through a policy-gated action boundary. Remote action, host commands, arbitrary filesystem access, mutation actions, provider-backed integrations, active learning promotion, and workflow mutation require stronger authority boundaries than bounded internal action execution.

\subsection{Autonomy Levels}

\begin{tabularx}{.95\linewidth}{@{}lX@{}}
	\toprule
	\textbf{Level} & \textbf{Meaning} \\
	\midrule
	Observation & Read local runtime state, existing knowledge, and safe metadata without mutation. \n
	Draft creation & Produce summaries, candidate records, or proposed changes without promoting them into active state. \n
	Bounded mutation & Modify authorized local state only when the action is logged, scoped, reversible where possible, and policy-permitted. \n
	External action & Contact remote systems, use external services, or disclose data only after classification, minimization, authorization, provenance, and recorded trace. \n
	Critical action & Credentials, identity, legal, financial, destructive, irreversible, or executable-code actions require a stronger human authorization path. \\
	\bottomrule
\end{tabularx}

\section{Evolution Control}

\subsection{Evolution Rules}

\begin{itemize}
	\item Runtime learning is represented by evaluated \texttt{LearningCandidate} proposals, not only by manual feedback collection.
	\item Learning strategies must declare what they are allowed to observe, propose, mutate, verify, and promote.
	\item A \texttt{LearningCandidate} cannot silently modify protected behavior.
	\item Learning proposals begin with source, evaluation evidence, scope, and rollback information before they are promoted.
	\item A \texttt{KnowledgeCandidate} is used only for pending knowledge-boundary change and requires provenance, truth metadata, correction paths, and recovery options.
	\item Saving a \texttt{KnowledgeRecord} is not learning by itself.
	\item Workflow and routing improvements may become default behavior through developer-defined checks, evaluation gates, and rollback policies, without requiring user authorization for every routine improvement.
	\item Model adaptation, fine-tuning, or weight updates are advanced model-evolution actions governed by stronger model-adaptation authority.
	\item Critical learning failures must be recoverable through snapshots, versioned state, VCS-backed records, or an equivalent rollback mechanism.
\end{itemize}

This policy keeps \Madre capable of evolving through conscious, checked learning without reducing improvement to constant user authorization. The final user should be interrupted as little as possible. Routine learning should be governed by developer-defined strategies, evidence, tests, promotion gates, \texttt{RuntimeJournal} traces, and rollback mechanisms; user authorization is reserved for boundaries where risk, privacy, identity, external action, or policy explicitly requires it.

D-011 fixes the controlled-promotion rule: \texttt{KnowledgeCandidate} denotes pending knowledge-boundary change and \texttt{LearningCandidate} denotes evaluated proposed module improvement. Neither candidate creation nor \texttt{KnowledgeRecord} storage silently rewrites protected behavior.

\madreChapter{Runtime Governance}
\section{Runtime Authority Rule}

Runtime authority belongs to the \Madre kernel and its policy boundaries. User turns, source text, generated material, action results, local inspection, \texttt{KnowledgeCandidate}, and \texttt{LearningCandidate} can inform decisions, but they cannot authorize their own authority, promote themselves, mutate protected state, or widen permissions. Bounded internal action use requires \texttt{PolicyDecision}, local access is read-first where it inspects state, and transitions are recorded through \texttt{RuntimeJournal}.

\section{Policy And Action Gate}

\begin{center}
\resizebox{.95\linewidth}{!}{
\begin{tikzpicture}[
	box/.style={draw, rounded corners=2pt, align=center, text width=3.3cm, minimum height=1.0cm, inner sep=5pt},
	gate/.style={draw, diamond, aspect=2, align=center, inner sep=2pt, fill=primaryColor!12},
	arrow/.style={-{Latex[length=2mm]}, thick}
]
	\node[box] (request) at (0,0) {Request, source,\\or generated material};
	\node[box] (classify) at (4,0) {Classify action,\\data, autonomy};
	\node[gate] (policy) at (8,0) {PolicyDecision};
	\node[box] (allow) at (12,1.8) {Allow bounded\\internal action};
	\node[box] (block) at (12,0) {Block, fail,\\or require authorization};
	\node[box] (candidate) at (12,-1.8) {KnowledgeCandidate\\or LearningCandidate};
	\node[box] (audit) at (16,0) {RuntimeJournal\\status, recovery trace};

	\draw[arrow] (request) -- (classify);
	\draw[arrow] (classify) -- (policy);
	\draw[arrow] (policy) -- (allow);
	\draw[arrow] (policy) -- (block);
	\draw[arrow] (policy) -- (candidate);
	\draw[arrow] (allow) -- (audit);
	\draw[arrow] (block) -- (audit);
	\draw[arrow] (candidate) -- (audit);
\end{tikzpicture}
}
\end{center}

\section{Policy Matrix}

\MADRETABLE5:{\textbf{Action class}&\textbf{Data class}& \textbf{Autonomy level} & \textbf{Authorization} & \textbf{Trace/Recovery}}{%
%\rowcolors{1}{secondaryColor!15!}{}
Read local metadata & Public or low-sensitivity local metadata & Observation & None by default & Log access when used for reasoning. \n
\midrule Draft text & Any allowed context bundle & Draft creation & None	unless sensitive scope is crossed & Retain as conversation or generated material, not authority or candidate by itself. \n
\midrule Update known material & Scoped module material & Bounded mutation & Required when provenance, truth metadata, scope, or sensitivity changes & Versioned \texttt{KnowledgeRecord} with correction and rollback; \texttt{KnowledgeCandidate} if boundary change is pending. \n
\midrule Invoke internal action & Small task-scoped non-sensitive input & Bounded internal action & Policy allow required & Requested action, input class, result, policy outcome, and recovery expectation recorded. \n
\midrule Contact remote service & Minimized and classified transfer set & External action & Explicit policy authorization & Transfer record, provider, purpose, and retention assumption recorded. \n
\midrule Modify credentials, identity, legal, financial, destructive, or executable state & Sensitive or critical data & Critical action & Strong human authorization	& Block unless recovery or explicit irreversible authorization exists. }

Knowledge-boundary updates, routing mutation, workflow mutation, and learning lifecycle behavior follow the same authority rule: a pending knowledge-boundary change is a \texttt{KnowledgeCandidate}; an evaluated behavior improvement is a \texttt{LearningCandidate}; neither becomes active merely by being produced.

\section{Learning Lifecycle}

Learning follows a controlled lifecycle: observation, evaluation, \texttt{LearningCandidate} creation, evidence collection, evaluation, promotion, active monitoring, and rollback. Knowledge handling remains distinct: a stored \texttt{KnowledgeRecord} is not learning, and a pending change to its owning knowledge boundary is a \texttt{KnowledgeCandidate}. Routing preferences, workflow improvements, and inference-selection improvements remain inactive unless the appropriate learning promotion rule validates them for a defined scope.

\madreChapter{Threat Model}

\section{Threat Taxonomy}

\MADRETABLE5:{\textbf{Threat} & \textbf{MADRE exposure} & \textbf{Primary control} &	\textbf{Verification surface}}{
	Source as instruction & Untrusted text attempts to alter policy, action execution, or knowledge handling by being interpreted as an instruction. & Source classification and instruction/data separation. &	Adversarial context scenario. \N
	Sensitive disclosure & Private context leaks to a model, action, or remote service. & Data classification, minimization, and remote-transfer authorization. &	Disclosure review. \N
	Excessive agency & generated material causes over-broad action. & Autonomy levels and action-class policy. & Policy matrix review. \N
	Unsafe action execution & Action invocation produces side effects outside task scope. &	Policy-gated internal action boundary; remote, host-command, filesystem, and mutation actions require stronger authority. & Allowed-action and blocked-action scenarios. \N
	Incorrect truth authority for known material & A \texttt{KnowledgeRecord} is used as factual authority despite unsuitable \texttt{truthLevel} or \texttt{truthAuthority}. & Required knowledge metadata, intended-use checks, correction, and reclassification path. & Non-authoritative knowledge-use and truth-correction scenario. \N
	Inspection misuse & Local console is treated as authority or mutation path.	& Read-first inspection boundary and shared policy model. & Local inspection negative scenario. \N
	Supply-chain or provider dependency & Model, dependency, or service changes behavior or availability. & Replaceable model backend and dependency policy. & Inference substitution review. \N
	Remote transfer risk & Local data is sent outside the user's domain without classification, minimization, authorization, or recorded trace. & Local-first default and remote-transfer policy boundary. & Remote-contact blocked scenario. \N
	Unbounded consumption & Delayed work consumes excessive compute, storage, or cost. & Resource budgets, queue limits, and status inspection. & Resource benchmark. \N
	Trace/recovery gaps & A result cannot be explained, corrected, invalidated, superseded, rolled back, or detached. & \texttt{RuntimeJournal}, append-only \texttt{RuntimeEvent}, and recovery semantics. & Recovery trace scenario. \N
	Inference output as false authority & Fluent generated material is treated as truth, policy, knowledge, behavior, routing, authority, or a candidate by itself. & Explicit generated-material handling and \texttt{PolicyDecision} boundary. & Output non-authority/non-candidate scenario.}

\madreChapter{Threat Catalogue}

\MADRETABLE3:{Threat family & \textbf{Primary boundary} & \textbf{Initial response}}{
	Source as instruction & Context bundle construction & Separate source material from instruction and policy. \N
	Sensitive disclosure & Data classification and remote transfer & Minimize, redact, authorize, and record disclosures. \N
	Unsafe action execution & Action and workflow contracts & Require action-class policy, bounded internal-action scope, and recovery expectation. \N
	Known material used under incorrect truth metadata & \texttt{KnowledgeRecord} use & Enforce provenance, \texttt{truthLevel}, \texttt{truthAuthority}, intended use, correction, and reclassification. \N
	Inspection bypass & Local console and status surfaces & Read-first inspection with no privileged mutation path. \N
	Generated-material false authority & generated material boundary & Treat generated material as produced material, not authority or candidate by itself. \N
	Provider dependency & Local Model Runtime Boundary & model backend substitution and local-first default. \N
	Unbounded consumption & Scheduler and resource controls & Queue, compute, storage, and cost limits.
}
